% Esquisse d'un programme (cote 119, copy 2): archive cote 119, site page 120 (PDF page 121), typescript page 42.
% typescript. Draft compared: esquisse/pages/esq_119_044.tex (copy 1, site 44).
% AI-transcribed by Claude Opus 5.5 (claude-opus-5-5, adaptive thinking, effort high), Message Batches API, batch esq119_full (msgbatch_015xqDHkTumDysRF1qngZhRJ), 2026-11-13; not yet checked by eye.
\scanpage{119/120}
\mspage{42}

correspondant doit être très étroite, du moins pour des $\mathcal{A}$ qui
peuvent être engendrés par un ``système de générateurs'' S fini -
mais il est à craindre que la moindre question intéressante qu'on
pourrait se poser sur cette situation soit actuellement hors de
portée ($^1$).

Ces réflexions anciennes ont repris quelque actualité
pour moi avec ma réflexion ultérieure sur les théories modérées.
Il me semble en effet qu'il est possible d'associer de façon na-%
turelle à une ``théorie différentiable'' $\mathcal{A}$ une théorie modérée $\mathcal{T}$
(ayant sans doute même corps de constantes), de telle façon que
tout $\mathcal{A}$-variété soit automatiquement munie d'une structure $\mathcal{T}$-modérée,
et inversement que pour tout espace $\mathcal{T}$-modéré compact $X$, on puisse
trouver une partie fermée modérée rare $Y$ dans $X$, telle que $X \smallsetminus Y$
provienne d'une $\mathcal{A}$-variété, et que de plus cette structure de $\mathcal{A}$-%
variété soit unique tout au moins dans le sens suivant : deux tel-%
les structures coïncident dans le complémentaire d'une partie modé-%
rée rare $Y' \supset Y$ de $X$. La théorie de dévissage des structures
modérées stratifiées (dont il a été question au par. précédent),
dans le cas de strates lisses, devrait d'ailleurs soulever des
questions beaucoup plus précises encore de comparaison des struc-%
tures modérées avec des structures de type différentiable (ou plu-%
tôt, $\mathbb{R}$-analytique). Je soupçonne que le type d'axiomatisation pro-%
posé ici pour la notion de ``théorie différentiable'' fournirait un
cadre naturel pour formuler de telles questions avec toute la pré-%
cision et toute la généralité souhaitables.


7. Depuis le mois de mars de l'an dernier, donc depuis près
d'un an, la plus grande partie de mon énergie a été consacrée à un
travail de réflexion sur les \uline{fondements de l'algèbre} (\uline{co})\uline{homologique}
\uline{non commutative}, ou ce qui revient au même, finalement, de l'\uline{algèbre}
\uline{homotopique}. Ces réflexions se sont concrétisées par un volumineux
paquet de notes dactylographiées, destinées à former le premier
volume (actuellement en cours d'achèvement) d'un ouvrage en deux

\margin{120}
